\documentclass[12pt]{article}
\usepackage[margin=1in]{geometry}
\usepackage[T1]{fontenc}

\begin{document}

\section*{Task 1: Requirements and Subsystems --- DocPilot}

\section{Functional and Non-Functional Requirements}

\subsection{Functional Requirements}

\begin{enumerate}
   \item \textbf{Grounded Document Q\&A}
   \begin{itemize}
      \item Users can query uploaded documents and receive answers with exact page and paragraph citations using a RAG pipeline.
      \item Architecturally significant because it requires tight integration of retrieval, context budgeting, and generation, forming the system's core pipeline.
   \end{itemize}

   \item \textbf{Scoped AI Editing with Diff View}
   \begin{itemize}
      \item Enables targeted edits with structured diff output and per-change accept/reject controls.
      \item Architecturally significant due to the need for Command + Memento patterns, ensuring reversibility and safe state transitions.
   \end{itemize}

   \item \textbf{Cross-Document Synthesis}
   \begin{itemize}
      \item Allows reasoning across multiple documents (comparison, contradiction detection, unified summaries).
      \item Significant because it requires multi-document retrieval, aggregation, and reasoning pipelines, increasing system complexity.
   \end{itemize}

   \item \textbf{RAG-Based Draft Generation}
   \begin{itemize}
      \item Generates structured outputs (reports, summaries) grounded in source documents with attribution.
      \item Architecturally important as it combines generation + citation tracking + context linking.
   \end{itemize}

   \item \textbf{Intelligent Model Routing and Context Sharing}
   \begin{itemize}
      \item Dynamically selects the best model based on task type, cost, and latency, while sharing context across models.
      \item Critical because it introduces a Model Gateway + Strategy pattern, central to system scalability and efficiency.
   \end{itemize}
\end{enumerate}

\subsection{Non-Functional Requirements}

\begin{enumerate}
   \item \textbf{Response Latency (NFR1)}
   \begin{itemize}
      \item Q\&A responses $\leq 10$ seconds; diff generation $\leq 5$ seconds.
      \item Significant as it constrains pipeline design, model selection, and caching strategies.
   \end{itemize}

   \item \textbf{Citation Accuracy (NFR2)}
   \begin{itemize}
      \item $\geq 80\%$ of answers must correctly cite source chunks.
      \item Drives design of retrieval precision, chunking strategy, and evaluation mechanisms.
   \end{itemize}

   \item \textbf{Retrieval Relevance (NFR3)}
   \begin{itemize}
      \item $\geq 85\%$ recall in top-3 retrieved chunks.
      \item Architecturally significant for vector database design and embedding quality.
   \end{itemize}

   \item \textbf{Parsing Fidelity (NFR4)}
   \begin{itemize}
      \item $\geq 95\%$ accuracy in extracting document structure (paragraphs, tables, etc.).
      \item Critical because all downstream tasks depend on correct parsing abstraction.
   \end{itemize}

   \item \textbf{Context Budget Management (NFR5)}
   \begin{itemize}
      \item Must not exceed $90\%$ of the model context window.
      \item Influences context engine design and prompt construction logic.
   \end{itemize}

   \item \textbf{Fallback Latency (NFR6)}
   \begin{itemize}
      \item Model fallback must start within 2 seconds.
      \item Requires Chain of Responsibility pattern and resilient orchestration.
   \end{itemize}
\end{enumerate}

\subsection{Architecturally Significant Requirements (Summary)}

The most critical requirements shaping the architecture are:

\begin{itemize}
   \item Grounded Q\&A and RAG pipeline $\rightarrow$ defines the core system workflow.
   \item Model routing and fallback $\rightarrow$ requires abstraction (Model Gateway, Strategy, Chain of Responsibility).
   \item Scoped editing with undo $\rightarrow$ enforces state management (Command + Memento).
   \item Cross-document synthesis $\rightarrow$ drives multi-agent orchestration.
   \item Latency and accuracy constraints $\rightarrow$ influence every layer (retrieval, parsing, generation).
\end{itemize}

\section{Subsystem Overview}

Based on the architecture described in the proposal (see diagram on page 2), the system can be decomposed into the following subsystems:

\subsection{Document Ingestion and Parsing Subsystem}
\begin{itemize}
   \item Handles document upload (PDF, DOCX, PPTX, MD).
   \item Uses adapters to convert different formats into a unified structure (block tree).
   \item Ensures high parsing fidelity for downstream processing.
\end{itemize}

\subsection{Context Engine (RAG Pipeline)}
\begin{itemize}
   \item Responsible for chunking, embedding, retrieval, and context budgeting.
   \item Core subsystem enabling grounded Q\&A and generation.
   \item Ensures relevant context is selected within model limits.
\end{itemize}

\subsection{Agent Orchestration Subsystem}
\begin{itemize}
   \item Coordinates specialized agents:
   \begin{itemize}
      \item Q\&A Agent
      \item Edit Agent
      \item Generate Agent
      \item Multi-Document Agent
   \end{itemize}
   \item Routes tasks and manages execution flow across pipeline stages.
\end{itemize}

\subsection{Model Gateway Subsystem}
\begin{itemize}
   \item Acts as a facade over multiple LLM providers (OpenAI, Anthropic, Google).
   \item Implements:
   \begin{itemize}
      \item Model selection (Strategy)
      \item Fallback handling (Chain of Responsibility)
   \end{itemize}
   \item Decouples the system from specific model APIs.
\end{itemize}

\subsection{Editing and Version Control Subsystem}
\begin{itemize}
   \item Implements Command + Memento patterns for reversible edits.
   \item Tracks changes and enables undo/redo functionality.
   \item Supports structured diff generation and user-controlled commits.
\end{itemize}

\subsection{Storage and Retrieval Subsystem}
\begin{itemize}
   \item Includes:
   \begin{itemize}
      \item Vector database (pgvector / Pinecone) for embeddings
      \item Metadata store for document structure
   \end{itemize}
   \item Uses Repository + CQRS pattern for separation of read/write paths.
\end{itemize}

\subsection{User Interface Subsystem}
\begin{itemize}
   \item Web-based interface (React + TypeScript).
   \item Provides:
   \begin{itemize}
      \item Document editor
      \item Assistant panel
      \item Multi-document workspace
   \end{itemize}
   \item Enables interaction with all system features.
\end{itemize}

\subsection{Event and Notification Subsystem}
\begin{itemize}
   \item Uses the Observer pattern to propagate document changes.
   \item Updates:
   \begin{itemize}
      \item Diff renderer
      \item Citation tracker
      \item Session memory
   \end{itemize}
   \item Ensures loose coupling between components.
\end{itemize}

\section*{Final Structure (Quick Summary)}

\textbf{Core Pipeline Subsystems}
\begin{itemize}
   \item Ingestion and Parsing
   \item Context Engine (RAG)
   \item Agent Orchestration
   \item Model Gateway
\end{itemize}

\textbf{Supporting Subsystems}
\begin{itemize}
   \item Editing and Version Control
   \item Storage and Retrieval
   \item UI Layer
   \item Event/Observer System
\end{itemize}

\end{document}

